% ------------------------------------------------------------
% CHAPTER 37
% ------------------------------------------------------------

\chapter{Toward a General Theory of Coherent Systems}

The preceding chapters have developed a progression of ideas moving from delta–sigma modulation as an engineering mechanism, through biological sensing and cognition, into the realm of semantic compression in artificial intelligence. Each stage revealed that the salient properties of delta–sigma modulation—oversampling, integration, collapse, entropy routing, error shaping, and adaptive feedback—appear wherever a system must maintain coherence while interacting with an uncertain world. These recurring structures suggest that delta–sigma modulation is not merely an algorithm or an engineering trick. Instead, it is a manifestation of a deeper mathematical pattern shared by all systems that must convert continuous fields into symbolic states while preserving identity under perturbation. This chapter develops that idea into a candidate for a general theory of coherent systems.

A coherent system may be defined as one that maintains a stable internal or external identity in the presence of noise, uncertainty, or external disturbance. This identity need not be conscious or biological; it may be mechanical, informational, cognitive, or semantic. The essential feature is that coherence persists across time, despite continual entropy injection. Persistence requires feedback loops that redistribute uncertainty so that entropy does not accumulate in regions where it would destroy the invariants defining the system’s identity. The delta–sigma modulator provides a minimal model: identity corresponds to the boundedness of the internal state and the shape of the noise spectrum; entropy enters through quantization; feedback routes this entropy into high-frequency channels; and coherence is maintained.

Generalizing this idea, consider a system described by a continuous field \(U(t)\) evolving on a manifold \(\mathcal{M}\). The field interacts with an environment that generates disturbances \(D(t)\). The system emits discrete or coarsely quantized symbols \(Y(t)\), which may represent actions, predictions, interpretations, or outputs more generally. These symbols enact collapse operations on the continuous state, discarding information and injecting entropy. The system possesses internal feedback mechanisms that update \(U(t)\) according to both its prior dynamics and the mismatch between its predicted and observed states. These mechanisms constitute the integrators, the loop filters, or the predictive machinery that maintain coherence. Coherence is defined by constraints on the evolution of \(U(t)\) that keep it within a stable region of \(\mathcal{M}\).

One may formalize these dynamics through a field equation of the form
\[
\frac{dU}{dt} = A(U) + B(D) - C(Y),
\]
where \(A\) encodes the system’s intrinsic dynamics, \(B\) represents external disturbance, and \(C\) represents the collapse induced by discrete symbolic actions. This equation generalizes the delta–sigma update to an arbitrary coherent system. The entropy injected by collapse acts as a forcing term. The system must then route this entropy out of the region of \(\mathcal{M}\) that defines its coherent identity. This routing is governed by the geometry of the manifold and the properties of the feedback law.

The geometry of coherence is central to this general theory. The manifold \(\mathcal{M}\) may possess a metric encoding the cost of deviations from the coherent region, a connection encoding how changes propagate, and a curvature encoding the susceptibility of the system to misalignment between internal predictions and external disturbances. Coherent systems must keep their trajectories close to geodesics within regions of low curvature. When disturbances push them into high-curvature regions, collapse events become more frequent or extreme, and feedback must compensate more vigorously. This geometric picture unifies the behavior of engineered delta–sigma modulators, neurons, cognitive systems, and adaptive AI architectures.

From the perspective of information geometry, the coherent region may be defined as a submanifold where divergence between predicted and observed fields remains bounded. The prediction error forms a vector in the tangent space of \(\mathcal{M}\), and the system must keep this vector within a cone of acceptable deviations. When the error grows too large, the system undergoes a collapse, projecting its representation back onto a stable symbolic manifold. These collapses correspond to quantization events in engineered systems, action potentials in neurons, conceptual shifts in cognition, and token emissions in AI models. At each collapse, the system discards information, injecting entropy that must be routed away from the coherent core.

The thermodynamic dimension of coherence emerges naturally. Systems maintain their invariants by expending energy to counteract entropy production. In delta–sigma modulators, energy is consumed by comparators, integrators, and switching transitions. In biological systems, metabolic energy supports ion transport, synaptic updates, and motor actions. In cognitive systems, the metaphorical energy corresponds to computational and attentional resources. In artificial intelligence, it corresponds to computational cost and gradient updates. The common pattern is that coherence requires work. Systems that cannot pay this cost lose coherence, diverge, or collapse into trivial or chaotic behavior. The conservation and expenditure of free energy thus become universal constraints.

Derived geometry provides the final layer of unification. When a system collapses a continuous field onto a discrete symbolic state, it performs a derived fiber product that generates hidden homotopies representing the mismatch between the continuous and discrete representations. Coherent systems must be capable of lifting these homotopies into corrected continuous states through feedback. Failure to perform these homotopy lifts leads to instability, overload, or fragmentation. In this sense, coherence is a homotopical property: a system is coherent to the extent that it can continually perform derived corrections to reconcile the differences between its continuous and symbolic representations.

This general framework suggests a universal architecture for coherent systems. At the base lies a continuous field representing the system’s state or beliefs. At the top lies a discrete symbolic layer representing actions, concepts, or emissions. Between them lie integrators or reservoirs that accumulate mismatch and feedback mechanisms that correct it. Entropy enters through collapse; coherence is maintained through routing; and stability is achieved through geometric alignment. Delta–sigma modulation, biological sensing, cognition, and artificial intelligence all instantiate variants of this architecture because it is the unique solution to the problem of coherence under measurement.

The general theory of coherent systems therefore positions delta–sigma modulation not as a special-purpose engineering method but as the archetypal form of coherence maintenance in dynamical systems. It unifies measurement, control, inference, and symbolic processing under a single mathematical umbrella. It reveals the deep structural constraints that govern all systems that maintain identity in the presence of noise. And it provides a roadmap for designing new coherent systems, from neuromorphic hardware to adaptive agents and distributed networks.

The next chapter brings these ideas into focus by presenting concluding observations, highlighting the implications for engineering, cognition, and physics, and synthesizing the unifying principle that measurement is the process of sculpting entropy to preserve coherent identities across time.

